% 乔治·基斯佳科夫斯基（George Kistiakowsky）（综述）
% license CCBYSA3
% type Wiki

本文根据 CC-BY-SA 协议转载翻译自维基百科\href{https://en.wikipedia.org/wiki/George_Kistiakowsky}{相关文章}。

\begin{figure}[ht]
\centering
\includegraphics[width=6cm]{./figures/a36c1f246ada9580.png}
\caption{乔治·基斯佳科夫斯基} \label{fig_QZJS_1}
\end{figure}
乔治·博格达诺维奇·基斯佳科夫斯基（俄语：Георгий Богданович Кистяковский；乌克兰语：Георгій Богданович Кістяківський，罗马化：Heorhii Bohdanovych Kistiakivskyi；1900年12月1日［旧历11月18日］—1982年12月7日）是一位乌克兰裔美国物理化学家，曾任哈佛大学教授，参与“曼哈顿计划”，并在后期担任美国总统德怀特·D·艾森豪威尔的科学顾问。

他出生于当时俄帝国境内的博亚尔卡（Boyarka）\(^\text{[1]}\)，出身于“一个古老的乌克兰哥萨克家族，该家族在俄国革命前是知识精英阶层的一部分”\(^\text{[2]}\)。在俄国内战期间，基斯佳科夫斯基逃离祖国，前往德国，在柏林大学师从马克斯·博登斯坦（Max Bodenstein），获得物理化学博士学位。1926年他移居美国，1930年加入哈佛大学任教，并于1933年成为美国公民。

第二次世界大战期间，基斯佳科夫斯基担任国家国防研究委员会（NDRC）爆炸物开发部门负责人，同时兼任爆炸物研究实验室（ERL）技术主任，负责监督包括RDX与HMX在内的新型炸药的研制。他还参与了爆炸的流体动力学理论研究以及成形装药技术的发展。1943年10月，他被聘为“曼哈顿计划”的顾问，不久后接管X部门，负责研发用于内爆型核武器的爆炸透镜。1945年7月，他亲眼见证了“三位一体”核试验中的首次原子爆炸。数周后，另一枚内爆型核武器“胖子”（Fat Man）被投放至长崎。

1962年至1965年间，基斯佳科夫斯基担任美国国家科学院科学、工程与公共政策委员会（COSEPUP）主席，并于1965年至1973年间出任该院副院长。由于反对越南战争，他切断了与政府的所有联系，并积极投身反战组织“宜居世界委员会”（Council for a Livable World），并于1977年出任该组织主席。
\subsection{早年生活}
乔治·博格达诺维奇·基斯佳科夫斯基于1900年12月1日［旧历11月18日］出生在俄国帝国基辅省的博亚尔卡（今属乌克兰）\(^\text{[1][2][3]}\)。他的祖父亚历山大·费多罗维奇·基斯佳科夫斯基是俄国帝国的法学教授兼律师，专攻刑法\(^\text{[5]}\)。其父博格丹·基斯佳科夫斯基是基辅大学的法理学教授\(^\text{[4]}\)，并于1919年当选为乌克兰国家科学院院士\(^\text{[6]}\)。母亲名为玛丽亚·别连德施塔姆（Maria Berendshtam），他还有一位哥哥亚历山大，后来成为鸟类学家\(^\text{[4]}\)。乔治的叔叔伊戈尔·基斯佳科夫斯基曾任乌克兰国政府内政部长\(^\text{[5]}\)。

基斯佳科夫斯基曾在基辅和莫斯科的私立学校就读，直到1917年俄国革命爆发。随后他加入了反共的白军。1920年，他乘坐一艘被征用的法国船只逃离俄国。此后，他辗转土耳其和南斯拉夫，最终前往德国，并于同年进入柏林大学学习\(^\text{[4]}\)。1925年，他在马克斯·博登斯坦（Max Bodenstein）的指导下获得物理化学博士学位，论文题目为“氯一氧化物与臭氧的光化学分解”。之后，他成为博登斯坦的研究助理\(^\text{[7][4]}\)。他发表的前两篇论文均基于博士论文的延伸研究，并与博登斯坦合作撰写\(^\text{[8]}\)。

1926年，基斯佳科夫斯基以国际教育委员会奖学金获得者的身份赴美。另一位博登斯坦的学生休·斯托特·泰勒（Hugh Stott Taylor）\(^\text{[9]}\)接受了导师对基斯佳科夫斯基的高度评价，为他提供了普林斯顿大学的研究岗位。同年，基斯佳科夫斯基与一位瑞典路德宗女性希尔德加德·梅比乌斯（Hildegard Moebius）结婚\(^\text{[4][10]}\)。1928年，他们的女儿维拉（Vera）出生。1972年，维拉成为麻省理工学院首位获聘物理学教授的女性\(^\text{[11]}\)。1927年，他的两年奖学金到期后，获得了研究助理职位与杜邦奖学金。1928年10月25日，他被任命为普林斯顿大学副教授(^7[4])。在此期间，他与泰勒共同发表了一系列论文\(^\text{[8]}\)，并在泰勒的鼓励下，出版了美国化学学会专著《光化学过程》(^8[4])。

1930年，基斯佳科夫斯基加入哈佛大学任教，并在此度过了整个学术生涯。在哈佛期间，他的研究领域包括热力学、光谱学和化学动力学。同时，他也越来越多地参与政府与工业界的技术咨询。1933年，他在哈佛大学晋升为副教授，并于同年成为美国公民，入选美国艺术与科学院\(^\text{[12]}\)。1938年，他被任命为“阿博特与詹姆斯·劳伦斯化学教授”\(^\text{[13]}\)。翌年，他当选为美国国家科学院院士\(^\text{[14]}\)。1940年，他又当选为美国哲学学会成员\(^\text{[15]}\)。
\subsection{第二次世界大战}
\subsubsection{国家国防研究委员会}
\begin{figure}[ht]
\centering
\includegraphics[width=6cm]{./figures/069a0bbe074c9f6a.png}
\caption{基斯佳科夫斯基在洛斯阿拉莫斯时期的战时安全证件} \label{fig_QZJS_2}
\end{figure}
预见到科学将在尚未参战的美国的二战中发挥更大作用，富兰克林·D·罗斯福总统于1940年6月27日成立了国家国防研究委员会（National Defense Research Committee，NDRC），由万尼瓦·布什（Vannevar Bush）出任主席。哈佛大学校长詹姆斯·B·康南特（James B. Conant）被任命为B部（负责炸弹、燃料、毒气与化学品）的负责人。他任命基斯佳科夫斯基负责A-1小组，A-1小组负责炸药研究。1941年6月，NDRC并入科学研究与发展办公室（Office of Scientific Research and Development，OSRD）。布什成为OSRD主席，康南特接替他出任NDRC主席，基斯佳科夫斯基则成为B部的负责人。1942年12月的一次重组中，B部被拆分，他出任第8部（负责炸药与推进剂）负责人，并一直担任该职直到1944年2月。

基斯佳科夫斯基对美国在炸药与推进剂方面的研究现状感到不满。康南特于1940年10月在宾夕法尼亚州布鲁塞顿（Bruceton）靠近矿务局实验室处设立了爆炸物研究实验室（Explosives Research Laboratory，ERL），基斯佳科夫斯基最初对其活动进行监督并偶有来访；但康南特直到1941年春才正式任命他为该实验室的技术主任。尽管起初受制于设施不足，ERL的规模从1941年的5名工作人员增长到1945年战时的最高峰——162名全职实验室人员。RDX是一项重要的研究方向。这种高能炸药在战前已由德国人开发，挑战在于建立一种可大规模工业生产的工艺。RDX还常与TNT混合制成B型炸药（Composition B），广泛用于各类弹药；以及用于鱼雷与深水炸弹的Torpex。试点工厂在1942年5月投入运行，1943年开始大规模生产。

为了响应一项特殊需求——为中国游击队研制一种能在日本检查站伪装通过的炸药，基斯佳科夫斯基将作为RDX工艺副产物的无毒炸药HMX与面粉混合，制成了名为“Aunt Jemima”（以一种煎饼粉品牌命名）的混合物。这是一种可食用的炸药，可以冒充普通面粉，甚至可以用来烹饪。

除了对RDX与HMX等合成炸药的研究外，爆炸物研究实验室还考察了爆轰与冲击波的性质。这一工作最初作为纯基础研究发起，未见明显或即时的应用。基斯佳科夫斯基于1941年和1942年访问英国，会见了包括威廉·潘尼（William Penney）与杰弗里·泰勒（Geoffrey Taylor）在内的英国专家。当基斯佳科夫斯基与埃德加·布莱特·威尔逊（Edgar Bright Wilson, Jr.）一起审视现有知识状况时，发现若干领域值得进一步研究。基斯佳科夫斯基开始着手研究查普曼—朱盖模型（Chapman–Jouguet model），该模型描述了爆轰产生的冲击波如何传播。

当时查普曼—朱盖模型的有效性仍有疑问，约翰·冯·诺依曼在普林斯顿高等研究院也在对此开展研究。基斯佳科夫斯基意识到，偏离流体动力学理论的原因在于化学反应本身的速率。要控制反应，需要把计算精确到微秒级别。第8部也被卷入了成形装药的研究，成形装药的机制由泰勒和詹姆斯·塔克（James Tuck）在1943年解释清楚。
\subsubsection{曼哈顿计划}
在曼哈顿计划的洛斯阿拉莫斯实验室，关于内爆装置的研究由塞思·内德梅耶（Seth Neddermeyer）领导，但他的部门主要用圆柱体与小装药进行试验，只得到过类似岩石的产物。由于普遍预计枪式核武器设计对铀-235和钚都可行且不需要内爆技术，他们的研究被列为低优先级。

1943年9月，洛斯阿拉莫斯实验室主任罗伯特·奥本海默邀请冯·诺依曼来访，用新的视角审查内爆研究。冯·诺依曼在审阅了内德梅耶的研究并与爱德华·泰勒讨论后，建议使用高能炸药成形装药使球体内爆，他指出这不仅可以比枪式方法更快地聚集裂变材料，而且能显著减少所需材料量。更高效的核武器设想给奥本海默、泰勒和汉斯·贝特留下了深刻印象，但他们认为需要一位炸药方面的专家。基斯佳科夫斯基的名字随即被提出，他于1943年10月作为顾问被引入该项目。
\begin{figure}[ht]
\centering
\includegraphics[width=6cm]{./figures/6d37fd5f5234999e.png}
\caption{在内爆型核武器中，多边形的炸药透镜围绕球形弹芯排列。} \label{fig_QZJS_3}
\end{figure}
基斯佳科夫斯基起初不愿前往，“部分原因”，他后来解释道，“我认为原子弹不会及时研制出来，而且我更希望帮助赢得战争”。在洛斯阿拉莫斯，他开始重组内爆研究工作。他引入了摄影和X光等技术来研究成形装药的行为。前者在爆炸物研究实验室（ERL）已被广泛使用，后者则在塔克的论文中有所描述，塔克还建议使用三维炸药透镜。正如曼哈顿计划的其他方面一样，炸药透镜的研究同时沿多条路线进行，因为正如基斯佳科夫斯基所指出的，“不可能预测哪些基本技术会更成功”。\(^\text{[31]}\)

基斯佳科夫斯基将对成形装药、Composition B等炸药以及ERL在1942–1943年使用的程序的详细知识带到了洛斯阿拉莫斯。ERL本身也逐渐被卷入内爆工作；其副主任邓肯·麦克杜格尔（Duncan MacDougall）也负责曼哈顿计划中的Q项目。1944年2月，基斯佳科夫斯基接替内德梅耶出任炸药（E）部主管。\(^\text{[30][32]}\)

在洛斯阿拉莫斯，埃米里奥·塞格雷（Emilio Segrè）小组证实反应堆中产出的钚含有过多的钚-240，使其不适用于枪式武器后，内爆计划获得了新的紧迫性。1944年7月的一系列危机会议得出结论：要想研制可用的钚武器，唯一的希望是内爆。8月，奥本海默对整个实验室进行了重组，集中力量攻关内爆。一支新的炸药小组——X部在基斯佳科夫斯基领导下成立，负责研制炸药透镜。\(^\text{[33][34]}\)

在基斯佳科夫斯基的领导下，X部设计出压缩裂变钚球所需的复杂炸药透镜。这些透镜采用两种爆速显著不同的炸药，以产生所需的波形。基斯佳科夫斯基为慢速炸药选择了巴拉托尔（Baratol）。在试验了多种速炸药后，X部最终选定了Composition B。将炸药铸造成正确形状的工作持续到1945年。透镜必须无瑕，必须研制出铸造Composition B和Baratol的工艺。ERL通过设计出一种便于铸造的Baratol制备程序而实现了这一目标。\(^\text{[35]}\)1945年3月，基斯佳科夫斯基成为所谓的“Cowpuncher委员会”的成员，该委员会负责协调内爆工作。\(^\text{[36]}\)1945年7月16日，基斯佳科夫斯基目睹了“三位一体”试验中首个装置的引爆。\(^\text{[37]}\)数周后，一枚内爆型武器“胖子”（Fat Man）被投放在长崎。\(^\text{[38]}\)

除内爆工作外，基斯佳科夫斯基还通过用炸药环砍伐树木为滑雪道造坡，促进了洛斯阿拉莫斯的滑雪活动——促成了Sawyer's Hill滑雪索道协会的成立。\(^\text{[39]}\)他于1942年与希尔德加德离婚，1945年与欧玛·E·舒勒（Irma E. Shuler）结婚，二人在1962年离婚，之后他与伊莱恩·马霍尼（Elaine Mahoney）结婚。\(^\text{[40]}\)
\subsection{白宫任职}
1957年，在艾森豪威尔政府时期，基斯佳科夫斯基被任命为总统科学顾问委员会成员，并于1959年接替詹姆斯·R·基利安（James R. Killian）出任主席。他于1959年至1961年间领导科学与技术政策办公室（Office of Science and Technology Policy），1961年由杰罗姆·B·维斯纳（Jerome B. Wiesner）接任。
\begin{figure}[ht]
\centering
\includegraphics[width=6cm]{./figures/9ce2f3467e266f3f.png}
\caption{“阿特拉斯”导弹，第一代洲际弹道导弹} \label{fig_QZJS_4}
\end{figure}
1958年，基斯佳科夫斯基向艾森豪威尔总统建议，仅依靠对外国军事设施的视察并不足以控制其核武器。他指出监控导弹潜艇的困难，并提出军备控制战略应重点放在裁军而非检查上。\(^\text{[42]}\)

1960年1月，作为军备控制规划与谈判的一部分，他提出了“阈值概念”（threshold concept）。根据这一提案，所有超出地震探测技术可检测范围的核试验都将被禁止。在此基础上，美国与苏联将共同努力改进探测技术，并随着技术进步逐步降低允许的试验当量。此举体现了所谓“国家技术验证手段”（national means of technical verification）——这是军控领域对情报收集活动的委婉称呼——可在不使现场检查要求高到苏联无法接受的情况下提供安全保障。美国于1960年1月在日内瓦军控会议上向苏联提出了这一概念，苏联在3月表示赞同，并建议以某一地震震级作为阈值。然而，随着1960年5月“U-2事件”的爆发，会谈最终破裂。\(^\text{[43]}\)

与此同时，在早期核军控工作进行的同时，参谋长联席会议主席内森·F·特怀宁（Nathan F. Twining）将军于1959年8月向国防部长尼尔·麦克埃尔罗伊（Neil McElroy）提交了一份备忘录\(^\text{[44]}\)，建议正式将国家核打击目标清单的制定以及统一核作战计划的准备工作交由战略空军司令部（SAC）负责。在此之前，陆、海、空三军各自独立制定目标计划，导致同一目标被多次打击，且各军种计划之间缺乏协调。例如，海军并未配合空军轰炸机的航线摧毁防空设施，导致作战效率低下。尽管特怀宁已将备忘录提交给麦克埃尔罗伊，但参谋长联席会议成员在1960年初仍未就该政策达成一致\(^\text{[45][46]}\)。继任国防部长托马斯·盖茨（Thomas Gates）随后请艾森豪威尔总统做出决定。\(^\text{[47]}\)

艾森豪威尔表示，他不会“把这种支离破碎、毫无协调的军事力量的怪物留给继任者”。1960年11月初，他派基斯佳科夫斯基前往内布拉斯加州奥马哈的战略空军司令部总部评估其作战计划。起初，基斯佳科夫斯基未获准查阅资料，艾森豪威尔随即下达更严厉的命令，要求SAC官员必须予以配合。基斯佳科夫斯基于11月29日提交的报告指出，现有计划之间缺乏协调，目标数量庞大，许多目标被多支部队重复打击，造成“过度杀伤”（overkill）。艾森豪威尔对这些计划感到震惊，随后不仅推动了单一综合作战计划（Single Integrated Operational Plan, SIOP）的制定，也重新审视了整个核作战的目标选择、需求生成与计划流程。\(^\text{[48]}\)
\subsection{晚年}
\begin{figure}[ht]
\centering
\includegraphics[width=6cm]{./figures/1d07cd2fad6c08e3.png}
\caption{阿瑟·M·波斯坎泽尔探访“基斯佳科夫斯基林”。} \label{fig_QZJS_5}
\end{figure}
在参与曼哈顿计划与白宫任职之间，以及离开白宫之后，基斯佳科夫斯基一直担任哈佛大学物理化学教授。1957年，当他被要求教授一年级课程时，他向自己印象深刻的讲师休伯特·阿利亚（Hubert Alyea）取经。阿利亚寄给他大约700张4×6英寸（10.2×15.2厘米）的卡片，上面详细记录了课堂实验的演示步骤。除了这些卡片外，基斯佳科夫斯基从不提前准备实验。他后来回忆道：

“我认为那样做是不给大自然女士一点发挥的机会。我走进教室，看一眼化学品和那堆卡片，然后对学生们说：‘让我们看看今天阿利亚为我们准备了什么。’我从不使用教材，只用你的卡片。我会扫一眼说明，直接做实验。如果成功了，我们就为你祈福，然后进行下一个实验；如果失败了，我们就咒骂你，并花剩下的整堂课试着让它成功。”\(^\text{[13]}\)

1972年，他从哈佛大学退休，被授予名誉教授头衔。\(^\text{[49]}\)

1962年至1965年间，基斯佳科夫斯基担任美国国家科学院科学、工程与公共政策委员会（COSEPUP）主席，1965年至1973年间任该院副院长。\(^\text{[50][51]}\)多年来他获得多项荣誉，包括1957年美国空军部颁发的“杰出平民服务奖章”（Decoration for Exceptional Civilian Service）。他还曾获杜鲁门总统颁发的“功绩奖章”（Medal for Merit）、艾森豪威尔总统于1961年颁发的“自由勋章”（Medal of Freedom），以及约翰逊总统于1967年颁发的“国家科学奖章”（National Medal of Science）。此外，他还获得美国化学学会颁发的1961年“查尔斯·拉思罗普·帕森斯公共服务奖”（Charles Lathrop Parsons Award）\(^\text{[52]}\)，1972年的“普里斯特利奖章”（Priestley Medal），以及哈佛大学授予的“富兰克林奖章”（Franklin Medal）\(^\text{[53][54]}\)。

晚年，基斯佳科夫斯基积极参与反战组织“宜居世界委员会”（Council for a Livable World），并因抗议美国卷入越南战争而与政府断绝关系。1977年，他出任该委员会主席，致力于反对核扩散。\(^\text{[49]}\)

1982年12月17日，基斯佳科夫斯基因癌症在马萨诸塞州剑桥市去世\(^\text{[55][54]}\)。其遗体火化，骨灰撒在马萨诸塞州科德角的夏季别墅附近\(^\text{[40]}\)。他的档案现保存在哈佛大学档案馆\(^\text{[56]}\)。此外，加利福尼亚州蒙哥马利森林州立自然保护区（Montgomery Woods State Natural Reserve）内的一片红杉林也被命名以纪念他。
\subsection{注释}
\begin{enumerate}
\item Кучеренко, Микола（2006）。“博亚尔卡的基斯佳科夫斯基故居”，《乌克兰档案学年鉴》（Український археографічний щорічник），第10/11期，第856页。
\item 《乌克兰评论》（The Ukrainian Review），第7卷（1959年），第125页。
\item 《1900年11月18日（旧历）乔治·基斯佳科夫斯基的出生记录》，载于博亚尔卡村圣母升天教堂的教区登记簿 // 乌克兰中央历史档案馆（ЦДІАК України），档号127，目录1078，案卷2220，页180背–181（俄语、乌克兰语）。
\item Dainton 1985，第379–380页。
\item Heuman 1998，第7–8页、第213页。
\item Heuman 1998，第37–38页。
\item “Kistiakowsky, George B. (George Bogdan), 1900–，《乔治·B·基斯佳科夫斯基档案目录》”，哈佛大学，存档于2020年2月12日，检索于2013年5月7日。
\item Dainton 1985，第401页。
\item “化学谱系——马克斯·欧内斯特·奥古斯特·博登斯坦学术家族树”，Chemistry Tree，检索于2013年5月7日。
\item “维拉·基斯佳科夫斯基访谈”，曼哈顿计划之声（Manhattan Project Voices）。
\item 《维拉·基斯佳科夫斯基文献集》（Vera Kistiakowsky Papers），MC 485，麻省理工学院档案与特藏馆，剑桥，马萨诸塞州，2019年7月1日存档。
\item “乔治·博格丹·基斯佳科夫斯基”，美国艺术与科学院（American Academy of Arts & Sciences），2023年2月9日，检索于2023年5月3日。
\item Dainton 1985，第382页。
\item “乔治·B·基斯佳科夫斯基”，美国国家科学院（nasonline.org），检索于2023年5月3日。
\item “美国哲学学会会员历史”，search.amphilsoc.org，检索于2023年5月3日。
\item Stewart 1948，第7页。
\item Stewart 1948，第10页。
\item Stewart 1948，第52–54页。
\item Stewart 1948，第88页。
\item Dainton 1985，第383页。
\item Noyes 1948，第25页。
\item Noyes 1948，第27页。
\item Noyes 1948，第28页。
\item Noyes 1948，第36–42页。
\item Clode, George（2012年6月12日）。“回到绘图板——‘Aunt Jemima’”，《军事历史月刊》（Military History Monthly），检索于2013年5月4日。
\item Noyes 1948，第51页。
\item Noyes 1948，第58–64页。
\item Chapman, David Leonard（1899年1月）。“论气体中的爆炸速率”，《哲学杂志》（Philosophical Magazine）第五辑，第47卷（284期），伦敦，第90–104页。doi:10.1080/14786449908621243。ISSN 1941-5982。LCCN sn86025845。
\item Noyes 1948，第73页。
\item Hoddeson 等（1993），第130–133页。

\end{enumerate}